% ==============================================================================
% 06_conclusion_outlook.tex – Chapter 6
% ==============================================================================
\chapter{Conclusion and Outlook}
\label{ch:fazit}

\section{Conclusion}
\label{sec:fazit}
This thesis studied how well a pretrained vision model can be adapted to one
person with a short webcam calibration. A ResNet-18 pretrained on ImageNet stayed
frozen, and only a small regression head was trained.

The most important results are:
\begin{itemize}
  \item A generic model trained on GazeCapture is not better than an estimate
        that always gives the screen centre. It therefore mainly serves as a
        starting point for personalization.
  \item With a realistic, temporal split, a few minutes of personal data lower
        the error from \SI{31.8}{\percent} to \SI{27.9}{\percent} of the
        screen diagonal (about $-$\SI{12}{\percent}, 18 sessions). In good
        sessions the gain stays large (about $-$\SI{46}{\percent}) and also
        holds for a new recording of the same person. One run takes about two
        seconds on a CPU.
  \item A random split gives \SI{19.8}{\percent} instead of
        \SI{30.3}{\percent} (about $-$\SI{35}{\percent}). It overestimates
        the benefit, however, because almost identical neighbouring frames end
        up in training and test.
  \item A calm background and good light are more important than the design of
        the calibration. A smooth movement of the target across the whole screen
        was best.
  \item Frames directly after a jump of the target have wrong labels. Removing
        them helps a little, but costs nothing.
  \item For a second person, personalization with the whole camera image brought
        no gain, but with the eye crop of the demo it did
        (\SI{26.9}{\percent} $\rightarrow$ \SI{9.7}{\percent}). The input is
        therefore decisive. A calibration is also tied to the person: the
        calibration of person~A was worse than the screen centre on person~B.
  \item A short moving average (about \SI{0.5}{\second}) makes the predictions
        about \SI{40}{\percent} calmer without a measurable loss of accuracy.
        Longer windows raise the error, because the prediction lags behind the
        gaze. Most of the error is systematic and cannot be removed by
        smoothing.
\end{itemize}
Overall, the thesis shows: data-efficient personalization works with simple
means. But the accuracy depends strongly on the recording conditions, and it must be
assessed with a realistic data split. The benefit also depends on the person
and on the input: for a second person, personalization only worked with an eye
crop, and a calibration of another person was worse than the screen centre. A
study with more persons has to show how general this is.

\section{Outlook}
\label{sec:ausblick}
Several next steps are obvious:
\begin{itemize}
  \item \textbf{Eye or face crops as input.} In the whole camera image, the eyes
        are very small. A crop gives the model much more gaze information; the
        demo run with person~B strongly supports this (\cref{sec:person_b}).
        This step should be the first to be tested systematically with all
        sessions.
  \item \textbf{Other pretrained models.} DINOv2, CLIP, MAE or UniGaze instead of
        ResNet-18, each with a frozen backbone and a small head.
  \item \textbf{Lightweight adaptation of the backbone.} Adapters or LoRA as a
        middle way between head-only training and full fine-tuning.
  \item \textbf{More persons and more sessions,} with conditions that change in
        only one factor at a time.
  \item \textbf{Adaptive smoothing.} Filters that smooth less during fast
        movement and more during fixations, for example the 1~Euro
        filter~\cite{casiez2012oneeuro}, could reach the calmness of the moving
        average without its lag.
  \item \textbf{Taking head movement into account,} for example with the head
        position as an extra feature or with normalized face images.
\end{itemize}

\paragraph{Web demo.}
As a practical test, I built a demo in the browser (server with FastAPI). The
user follows a point for about \SI{60}{\second}. The point starts at a random
place, moves in a random direction at constant speed and bounces off the
edges. A path is only used if it visits all cells of a $3\times3$ grid; this
keeps the full screen coverage that proved important in
\cref{sec:session_design}, while the calibration feels less monotonous. Then a
model is personalized, and the gaze is shown live as a light field on a grid of
$3\times3$ cells. The displayed position is a causal moving average over the
last \SI{0.5}{\second} (\cref{sec:smoothing}). At the end, an evaluation shows
when which cell was looked at. The demo already uses some ideas from the
outlook: an eye crop, spatial features of the ResNet-18, a ridge regression and
the head position as an extra feature. A short test with five points
(\SI{3}{\second} each) after the calibration measures the error on new frames.
In two runs of the new version (person~A), the error on the
validation data of the calibration was \SI{19.1}{\percent} and
\SI{6.7}{\percent}; the base model without adaptation reached
\SI{25.4}{\percent} and \SI{26.0}{\percent}, and the constant estimate
\SI{28.3}{\percent} in the second run. In the second run, the five-point test
reached \SI{11.8}{\percent} (78 measurements); the best method there was the
ridge regression on spatial features, and training took \SI{6.1}{\second}. In
both runs, the base model pointed only at the centre cell during the whole
tracking, while the personalized model spread the gaze over eight and all nine
cells. A run with person~B reached \SI{9.7}{\percent} on the validation data
and \SI{20.8}{\percent} in the five-point test (\cref{tab:demo_persons}). With the earlier version of the demo (fixed snake
path), the error in the five-point test was between about \SI{17.6}{\percent} and
\SI{25.2}{\percent} in several single tests (last test: \SI{19.0}{\percent});
a constant estimate reaches about \SI{23}{\percent} on the calibration data. In
a direct comparison, the test error fell from \SI{23.0}{\percent} to
\SI{17.6}{\percent} with the head position. The values vary strongly between
runs and were not evaluated systematically. But they show that head movement
and light are large sources of error in live use.
